


\subsection{Reweighting Design Exploration}







\begin{algorithm}[H]
    \caption{Reinforced Flow Matching}
    \label{alg:velocity-add-delete}
    \textbf{Require:} Initial model $\theta_0$, reward function $R$, prompt distribution $p(c)$, number of prompts per iteration $C$, rollout size per prompt $K$, batch size $B$, total samples per iteration $D = KC$, rollout coverage coefficient $\beta$, mass-shift clip $\Delta_{\mathrm{clip}}$, reference regularization $\lambda_{\mathrm{ref}}$, EMA rate $\eta$, reward-scale constant $\epsilon_R$, numerical constant $\epsilon_{\mathrm{alg}}$, gradient accumulation steps $N_{\mathrm{acc}}$.\\
    \textbf{Initialize:} Training policy $\theta \leftarrow \theta_0$, old policy $\theta_{\mathrm{old}} \leftarrow \theta_0$, reference policy $\theta_{\mathrm{ref}} \leftarrow \theta_0$, data buffer $\mathcal{D} \leftarrow \emptyset$.
    \begin{algorithmic}[1]
        \For{\text{each iteration}}
            \For{\text{each prompt $c_j \sim p(c)$, $j = 1, \dots, C$}} \Comment{Rollout Step, Data Collection}
                \State Sample $K$ clean images $\{\mathbf{x}_0^{j,k}\}_{k=1}^{K}$ with $\theta$ conditioned on $c_j$.
                \State Evaluate rewards $R_{j,k} \leftarrow R(\mathbf{x}_0^{j,k}, c_j)$ for $k = 1, \dots, K$.
            \EndFor

            \State Compute prompt-wise reward means
            $\bar{R}_{c_j} \leftarrow \frac{1}{K}\sum_{k=1}^{K} R_{j,k}$ for $j = 1, \dots, C$.

            \State Compute global reward mean
            $\bar{R} \leftarrow \frac{1}{D}\sum_{j=1}^{C}\sum_{k=1}^{K} R_{j,k}$.

            \State Compute global reward scale
            $\sigma_R \leftarrow \sqrt{\frac{1}{D}\sum_{j=1}^{C}\sum_{k=1}^{K}(R_{j,k} - \bar{R})^2}.$

            \For{\text{each prompt $c_j$, $j = 1, \dots, C$}}
                \State Compute normalized weights
                \[
                    \Delta_{j,k} \leftarrow \frac{1}{\Delta_{\mathrm{clip}}}
                    \operatorname{clip}\!\left(
                        \frac{R_{j,k} - \bar{R}_{c_j}}{\sigma_R},
                        -\Delta_{\mathrm{clip}},
                        \Delta_{\mathrm{clip}}
                    \right),
                    \quad
                    w_{j,k}-1\leftarrow \Delta_{j,k},
                    \quad k = 1, \dots, K.
                \]

                \State Store
                $\mathcal{D}\leftarrow\mathcal{D}\cup
                \{(c_j,\mathbf{x}_0^{j,k},w_{j,k})\}_{k=1}^{K}$.
            \EndFor

            \For{\text{each mini batch }
            $\{(c_i,\mathbf{x}_i^0,w_i)\}_{i=1}^{B}\subset\mathcal{D}$}
                \State Sample $t_i\sim\mathcal{U}(0,1)$,
                $\boldsymbol{\epsilon}_i\sim\mathcal{N}(\mathbf{0},\mathbf{I})$;
                set $\mathbf{x}_i^{t_i}\leftarrow
                (1-t_i)\mathbf{x}_i^0+t_i\boldsymbol{\epsilon}_i$.
                \State Predict $\mathbf{x}_i^\theta$,
                $\mathbf{x}_i^{\mathrm{old}}$, and $\mathbf{x}_i^{\mathrm{ref}}$
                from $(\mathbf{x}_i^{t_i},c_i,t_i)$ with
                $\theta$, $\theta_{\mathrm{old}}$, and $\theta_{\mathrm{ref}}$.
                \State $\mathbf{x}_i^{\mathrm{target}}\leftarrow
                \mathbf{x}_i^{\mathrm{old}}+
                \beta(w_i-1)(\mathbf{x}_i^0-\mathbf{x}_i^{\mathrm{old}})$.
                \State Compute
                \[
                    \begin{aligned}
                    \mathcal{L}_{\mathrm{policy}}
                    &\leftarrow\frac{1}{B}\sum_{i=1}^{B}
                    \frac{\|\mathbf{x}_i^\theta-\mathbf{x}_i^{\mathrm{target}}\|_2^2}
                    {\operatorname{sg}[\operatorname{mean}
                    |\mathbf{x}_i^0-\mathbf{x}_i^{\mathrm{old}}|]},\\
                    \mathcal{L}_{\mathrm{ref}}
                    &\leftarrow\frac{1}{B}\sum_{i=1}^{B}
                    \|\mathbf{v}_i^\theta-\mathbf{v}_i^{\mathrm{ref}}\|_2^2.
                    \end{aligned}
                \]
                \State Backpropagate
                $(\Delta_{clip}\mathcal{L}_{\mathrm{policy}}+
                \lambda_{\mathrm{ref}}\mathcal{L}_{\mathrm{ref}})/N_{\mathrm{acc}}$.
            \EndFor
            \State Update $\theta$; set
            $\theta_{\mathrm{old}}\leftarrow
            \eta\theta_{\mathrm{old}}+(1-\eta)\theta$ and
            $\mathcal{D}\leftarrow\emptyset$.
        \EndFor
    \end{algorithmic}
    \textbf{Output:} $\mathbf{v}_\theta$
\end{algorithm}


\begin{algorithm}[H]
    \scriptsize
    \setlength{\abovedisplayskip}{1pt}
    \setlength{\belowdisplayskip}{1pt}
    \renewcommand{\arraystretch}{0.9}
    \caption{Scheme B: Symmetric Hard Selection}
    \label{alg:x-prediction-symmetric-hard-selection}
    \textbf{Require:} Initial model $\theta_0$, reward function $R$,
    prompt distribution $p(c)$, $C$
    prompts per iteration, $K$ samples per prompt, $D=CK$, batch size $B$,
    retained ratio $\rho$, shift strength $\beta$, mass-shift clip
    $\Delta_{\mathrm{clip}}$, reference regularization $\lambda_{\mathrm{ref}}$,
    EMA rate $\eta$, accumulation steps $N_{\mathrm{acc}}$.\\
    \textbf{Initialize:}
    $\theta,\theta_{\mathrm{old}},\theta_{\mathrm{ref}}\leftarrow\theta_0$,
    $\mathcal{D}\leftarrow\emptyset$.
    \begin{algorithmic}[1]
        \For{\text{each iteration}}
            \For{\text{each prompt $c_j\sim p(c)$, $j=1,\dots,C$}}
                \State Sample $\{\mathbf{x}_0^{j,k}\}_{k=1}^{K}$ with $\theta$
                conditioned on $c_j$ and set
                $R_{j,k}\leftarrow R(\mathbf{x}_0^{j,k},c_j)$.

            \EndFor
            \State Compute global reward mean
            $\bar{R}\leftarrow \frac{1}{D}\sum_{j=1}^{C}\sum_{k=1}^{K}R_{j,k}$.
            \State Compute global reward scale
            $\sigma_R \leftarrow \sqrt{\frac{1}{D}\sum_{j=1}^{C}\sum_{k=1}^{K}(R_{j,k} - \bar{R})^2}.$

            \For{\text{each prompt $c_j\sim p(c)$, $j=1,\dots,C$}}

                \State Set the mass-shift center to the prompt-group mean:
                $b_j\leftarrow K^{-1}\sum_{k=1}^{K}R_{j,k}$.
                \State Within this prompt group, compute
                \[
                    % a_j^-\leftarrow\max_{R_{j,k}<b_j}|R_{j,k}-b_j|,
                    % \quad
                    % a_j^+\leftarrow\max_{R_{j,k}>b_j}|R_{j,k}-b_j|,
                    % \quad
                    \Delta_{j,k}\leftarrow
                    \frac{1}{\Delta_{\mathrm{clip}}}
                    \operatorname{clip}\!\left(
                        \frac{R_{j,k}-b_j}{\sigma_R},
                        -\Delta_{\mathrm{clip}},
                        \Delta_{\mathrm{clip}}
                    \right)
                \]
                \State Sort the $K$ samples so that
                $\Delta_{j,1}\le\cdots\le\Delta_{j,q_j}\le0
                \le\Delta_{j,q_j+1}\le\cdots\le\Delta_{j,K}$, and set
                \[
                    M_j^-\leftarrow\left|\sum_{k=1}^{q_j}\Delta_{j,k}\right|,
                    \qquad
                    M_j^+\leftarrow\sum_{k=q_j+1}^{K}\Delta_{j,k}.
                \]
                \State Select the two strong-tail boundaries:
                \[
                    \begin{aligned}
                        l_{j,\rho}&\leftarrow
                        \arg\min_{1\le l\le q_j}
                        \left|\left|\sum_{k=1}^{l}\Delta_{j,k}\right|
                        -\rho M_j^-\right|,\\
                        h_{j,\rho}&\leftarrow
                        \arg\min_{q_j+1\le h\le K}
                        \left|\sum_{k=h}^{K}\Delta_{j,k}
                        -\rho M_j^+\right|.
                    \end{aligned}
                \]
                \State Conduct selection for both sides:
                \[
                    s_{j,k}\leftarrow
                    \begin{cases}
                        \Delta_{j,k},&k\le l_{j,\rho},\\
                        0,&l_{j,\rho}<k<h_{j,\rho},\\
                        \Delta_{j,k},&k\ge h_{j,\rho}.
                    \end{cases}
                \]
                \State Normalize the selected mass within this prompt group:
                \[
                    % m_j^-\leftarrow\sum_{s_{j,k}<0}(-s_{j,k}),\quad
                    % m_j^+\leftarrow\sum_{s_{j,k}>0}s_{j,k},\quad
                    % w_{j,k}-1\leftarrow
                    % \begin{cases}
                    %     s_{j,k}/m_j^-,&s_{j,k}<0,\\
                    %     0,&s_{j,k}=0,\\
                    %     s_{j,k}/m_j^+,&s_{j,k}>0.
                    % \end{cases}
                    w_{j,k}-1\leftarrow s_{j,k}/\rho
                \]
                \State Store
                $\mathcal{D}\leftarrow\mathcal{D}\cup
                \{(c_j,\mathbf{x}_0^{j,k},w_{j,k})\}_{k=1}^{K}$.
            \EndFor

            \For{\text{each mini batch }
            $\{(c_i,\mathbf{x}_i^0,w_i)\}_{i=1}^{B}\subset\mathcal{D}$}
                \State Sample $t_i\sim\mathcal{U}(0,1)$,
                $\boldsymbol{\epsilon}_i\sim\mathcal{N}(\mathbf{0},\mathbf{I})$;
                set $\mathbf{x}_i^{t_i}\leftarrow
                (1-t_i)\mathbf{x}_i^0+t_i\boldsymbol{\epsilon}_i$.
                \State Predict $\mathbf{x}_i^\theta$,
                $\mathbf{x}_i^{\mathrm{old}}$, and $\mathbf{x}_i^{\mathrm{ref}}$
                from $(\mathbf{x}_i^{t_i},c_i,t_i)$ with
                $\theta$, $\theta_{\mathrm{old}}$, and $\theta_{\mathrm{ref}}$.
                \State $\mathbf{x}_i^{\mathrm{target}}\leftarrow
                \mathbf{x}_i^{\mathrm{old}}+
                \beta(w_i-1)(\mathbf{x}_i^0-\mathbf{x}_i^{\mathrm{old}})$.
                \State Compute
                \[
                    \begin{aligned}
                    \mathcal{L}_{\mathrm{policy}}
                    &\leftarrow\frac{1}{B}\sum_{i=1}^{B}
                    \frac{\|\mathbf{x}_i^\theta-\mathbf{x}_i^{\mathrm{target}}\|_2^2}
                    {\operatorname{sg}[\operatorname{mean}
                    |\mathbf{x}_i^0-\mathbf{x}_i^{\mathrm{old}}|]},\\
                    \mathcal{L}_{\mathrm{ref}}
                    &\leftarrow\frac{1}{B}\sum_{i=1}^{B}
                    \|\mathbf{v}_i^\theta-\mathbf{v}_i^{\mathrm{ref}}\|_2^2.
                    \end{aligned}
                \]
                \State Backpropagate
                $(\Delta_{\mathrm{clip}}\mathcal{L}_{\mathrm{policy}}+
                \lambda_{\mathrm{ref}}\mathcal{L}_{\mathrm{ref}})/N_{\mathrm{acc}}$.
            \EndFor
            \State Update $\theta$; set
            $\theta_{\mathrm{old}}\leftarrow
            \eta\theta_{\mathrm{old}}+(1-\eta)\theta$ and
            $\mathcal{D}\leftarrow\emptyset$.
        \EndFor
    \end{algorithmic}
    \textbf{Output:} $\mathbf{x}_\theta$
\end{algorithm}


\begin{algorithm}[H]
    \scriptsize
    \setlength{\abovedisplayskip}{1pt}
    \setlength{\belowdisplayskip}{1pt}
    \renewcommand{\arraystretch}{0.9}
    \caption{Scheme C: Asymmetric Hard Selection}
    \label{alg:x-prediction-asymmetric-hard-selection}
    \textbf{Require:} Initial model $\theta_0$, reward function $R$,
    prompt distribution $p(c)$, $C$ prompts per iteration, $K$ samples per
    prompt, $D=CK$, batch size $B$, retained ratio $\rho$, shift strength
    $\beta$, mass-shift clip $\Delta_{\mathrm{clip}}$, reference coefficient
    $\lambda_{\mathrm{ref}}$, EMA rate $\eta$,
    accumulation steps $N_{\mathrm{acc}}$.\\
    \textbf{Initialize:}
    $\theta,\theta_{\mathrm{old}},\theta_{\mathrm{ref}}\leftarrow\theta_0$,
    $\mathcal{D}\leftarrow\emptyset$.
    \begin{algorithmic}[1]
        \For{\text{each iteration}}
            \For{\text{each prompt $c_j\sim p(c)$, $j=1,\dots,C$}}
                \State Sample $\{\mathbf{x}_0^{j,k}\}_{k=1}^{K}$ with $\theta$
                conditioned on $c_j$ and set
                $R_{j,k}\leftarrow R(\mathbf{x}_0^{j,k},c_j)$.

            \EndFor

            \State Compute global reward mean
            $\bar{R}\leftarrow \frac{1}{D}\sum_{j=1}^{C}\sum_{k=1}^{K}R_{j,k}$.
            \State Compute global reward scale
            $\sigma_R \leftarrow \sqrt{\frac{1}{D}\sum_{j=1}^{C}\sum_{k=1}^{K}(R_{j,k} - \bar{R})^2}.$

            \For{\text{each prompt $c_j\sim p(c)$, $j=1,\dots,C$}}

                \State Set the mass-shift center to the prompt-group mean:
                $b_j\leftarrow K^{-1}\sum_{k=1}^{K}R_{j,k}$.
                \State Within this prompt group, compute
                \[
                    % a_j^-\leftarrow\max_{R_{j,k}<b_j}|R_{j,k}-b_j|,
                    % \quad
                    % a_j^+\leftarrow\max_{R_{j,k}>b_j}|R_{j,k}-b_j|,
                    % \quad
                    \Delta_{j,k}\leftarrow
                    \frac{1}{\Delta_{\mathrm{clip}}}
                    \operatorname{clip}\!\left(
                        \frac{R_{j,k}-b_j}{\sigma_R},
                        -\Delta_{\mathrm{clip}},
                        \Delta_{\mathrm{clip}}
                    \right)
                \]
                \State Sort the $K$ samples so that
                $\Delta_{j,1}\le\cdots\le\Delta_{j,q_j}\le0
                \le\Delta_{j,q_j+1}\le\cdots\le\Delta_{j,K}$.
                \State Select the positive strong-tail boundary:
                \[
                    M_j^+\leftarrow\sum_{k=q_j+1}^{K}\Delta_{j,k},
                    \qquad
                    h_{j,\rho}\leftarrow
                    \arg\min_{q_j+1\le h\le K}
                    \left|\sum_{k=h}^{K}\Delta_{j,k}-\rho M_j^+\right|.
                \]
                \State Keep all negative samples and the positive strong signal samples:
                \[
                    s_{j,k}\leftarrow
                    \begin{cases}
                        \Delta_{j,k},&k\le q_j,\\
                        0,&q_j<k<h_{j,\rho},\\
                        \Delta_{j,k},&k\ge h_{j,\rho}.
                    \end{cases}
                \]
                \State Normalize the selected mass within this prompt group:
                \[
                    % m_j^-\leftarrow\sum_{s_{j,k}<0}(-s_{j,k}),\quad
                    % m_j^+\leftarrow\sum_{s_{j,k}>0}s_{j,k},\quad
                    % w_{j,k}-1\leftarrow
                    % \begin{cases}
                    %     s_{j,k}/m_j^-,&s_{j,k}<0,\\
                    %     0,&s_{j,k}=0,\\
                    %     s_{j,k}/m_j^+,&s_{j,k}>0.
                    % \end{cases}
                    w_{j,k}-1\leftarrow
                    \begin{cases}
                        s_{j,k},&s_{j,k}<0,\\
                        0,&s_{j,k}=0,\\
                        s_{j,k}/\rho,&s_{j,k}>0.
                    \end{cases}                    
                \]
                \State Store
                $\mathcal{D}\leftarrow\mathcal{D}\cup
                \{(c_j,\mathbf{x}_0^{j,k},w_{j,k})\}_{k=1}^{K}$.
            \EndFor

            \For{\text{each mini batch }
            $\{(c_i,\mathbf{x}_i^0,w_i)\}_{i=1}^{B}\subset\mathcal{D}$}
                \State Sample $t_i\sim\mathcal{U}(0,1)$,
                $\boldsymbol{\epsilon}_i\sim\mathcal{N}(\mathbf{0},\mathbf{I})$;
                set $\mathbf{x}_i^{t_i}\leftarrow
                (1-t_i)\mathbf{x}_i^0+t_i\boldsymbol{\epsilon}_i$.
                \State Predict $\mathbf{x}_i^\theta$,
                $\mathbf{x}_i^{\mathrm{old}}$, and $\mathbf{x}_i^{\mathrm{ref}}$
                from $(\mathbf{x}_i^{t_i},c_i,t_i)$ with
                $\theta$, $\theta_{\mathrm{old}}$, and $\theta_{\mathrm{ref}}$.
                \State $\mathbf{x}_i^{\mathrm{target}}\leftarrow
                \mathbf{x}_i^{\mathrm{old}}+
                \beta(w_i-1)(\mathbf{x}_i^0-\mathbf{x}_i^{\mathrm{old}})$.
                \State Compute
                \[
                    \begin{aligned}
                    \mathcal{L}_{\mathrm{policy}}
                    &\leftarrow\frac{1}{B}\sum_{i=1}^{B}
                    \frac{\|\mathbf{x}_i^\theta-\mathbf{x}_i^{\mathrm{target}}\|_2^2}
                    {\operatorname{sg}[\operatorname{mean}
                    |\mathbf{x}_i^0-\mathbf{x}_i^{\mathrm{old}}|]},\\
                    \mathcal{L}_{\mathrm{ref}}
                    &\leftarrow\frac{1}{B}\sum_{i=1}^{B}
                    \|\mathbf{v}_i^\theta-\mathbf{v}_i^{\mathrm{ref}}\|_2^2.
                    \end{aligned}
                \]
                \State Backpropagate
                $(\Delta_{\mathrm{clip}}\mathcal{L}_{\mathrm{policy}}+
                \lambda_{\mathrm{ref}}\mathcal{L}_{\mathrm{ref}})/N_{\mathrm{acc}}$.
            \EndFor
            \State Update $\theta$; set
            $\theta_{\mathrm{old}}\leftarrow
            \eta\theta_{\mathrm{old}}+(1-\eta)\theta$ and
            $\mathcal{D}\leftarrow\emptyset$.
        \EndFor
    \end{algorithmic}
    \textbf{Output:} $\mathbf{x}_\theta$
\end{algorithm}



Schemes B and C use the same normalized and clipped signal
$\Delta_{j,k}$ as the base scheme.  Scheme B selects strong-tail samples on
both sides, whereas Scheme C keeps all negative samples and selects only the
positive strong tail.

The selected mass is approximately a fraction $\rho$ of the corresponding
base mass.  Dividing the retained signal by $\rho$ therefore approximately
preserves the total mass of the base scheme while reallocating it from
borderline samples to strong samples.  In Scheme C, the negative side is not
selected and is consequently not divided by $\rho$.

When $\rho=1$, no nonzero sample is removed and both schemes satisfy
$w_{j,k}-1=\Delta_{j,k}$, so they reduce exactly to the base scheme under the
same training hyperparameters.  When $\rho<1$, the retained samples receive
$1/\rho$ times their base signal and the discarded samples receive zero; the
total mass is only approximately preserved because the selection boundary
must fall between discrete samples.  Thus these schemes mainly test the
effect of concentrating the base training signal on stronger samples, although
a small $\rho$ may produce large individual target shifts.


\subsection{Mass-Shift Experiment Matrix}



\begin{table}[H]
    \centering
    \caption{Settings for the nineteen mass-shift experiments.  The launcher
    names below share the prefix
    \texttt{run\_sd3\_4090\_16gpu\_nft-ours-4.selection-}.}
    \label{tab:mass-shift-hard-selection-sweep}
    \small
    \setlength{\tabcolsep}{4pt}
    \renewcommand{\arraystretch}{1.15}
    \begin{tabularx}{\textwidth}{@{}X c c c c@{}}
        \toprule
        Experiment group & Scheme & $\rho$ & $\beta$ & Runs \\
        \midrule
        Equivalence baseline & B ($=$ C) & $1$ & $1.0$ & $1$ \\
        \midrule
        Symmetric hard selection & B
        & $\{1/3,1/2,2/3\}$ & $\{0.2,1.0,5.0\}$ & $9$ \\
        \midrule
        Asymmetric hard selection & C
        & $\{1/3,1/2,2/3\}$ & $\{0.2,1.0,5.0\}$ & $9$ \\
        \bottomrule
    \end{tabularx}

    \vspace{0.4em}
    \begin{minipage}{0.98\textwidth}
        \footnotesize
        The baseline launcher suffix is
        \texttt{baseline-rho1-beta1.0.sh}.  For B and C, each listed
        $\rho$ is crossed with all three $\beta$ values.  The $\beta=1.0$
        launchers retain their existing suffixes, for example
        \texttt{b-rho-1of3.sh}; the added launchers append
        \texttt{-beta0.2} or \texttt{-beta5.0}, for example
        \texttt{b-rho-1of3-beta0.2.sh}.
    \end{minipage}

    \vspace{0.6em}
    \begin{tabularx}{\textwidth}{@{}>{\bfseries}l X@{}}
        \toprule
        Shared item & Setting used by all nineteen experiments \\
        \midrule
        Model and task
        & SD 3.5 Medium, GenEval reward evaluated with
          OpenCLIP, $512\times512$ resolution. \\
        Rollout
        & $K=24$ images per prompt, 48 prompt groups (1152 images) per
          iteration, 10 deterministic DPM2 sampling steps, guidance scale
          $1.0$. \\
        Compute and batch
        & 16 RTX 4090 GPUs ($2$ nodes $\times$ $8$ GPUs), per-device batch
          size $6$, global micro-batch $96$, gradient accumulation $12$, and
          effective batch size $1152$. \\
        Objective
        & Advantage clip $\Delta_{\mathrm{clip}}=5$,
          $\beta=\texttt{config.beta}$ as listed in each row,
          and reference/KL
          coefficient $\lambda_{\mathrm{ref}}=\texttt{config.train.beta}=10^{-4}$. \\
        Training
        & Full-timestep training
          (\texttt{timestep\_fraction}$=1.0$),
          \texttt{trajectory\_alpha\_prediction=old\_prediction}, learning rate $3\times10^{-4}$, and maximum gradient norm
          $1.0$. \\
        \bottomrule
    \end{tabularx}
\end{table}

\paragraph{Baseline equivalence check.}
The run with $\rho=1$ and $\beta=1$ verifies that the selection
implementation reduces to the base scheme: no nonzero sample is discarded,
the $1/\rho$ rescaling is the identity, and hence
$w_{j,k}-1=\Delta_{j,k}$.  The expected result is that its learning curve and
final metrics match the base scheme up to ordinary run-to-run variation.

\paragraph{Scheme B: symmetric hard selection.}
The $3\times3$ sweep separates the effect of retained mass $\rho$ from the
overall shift strength $\beta$.  Comparisons at fixed $\beta$ test whether
concentrating both positive and negative signals on strong-tail samples is
better than retaining borderline samples; comparisons at fixed $\rho$ test
the required target-shift magnitude.  The intended conclusion is whether a
moderate $\rho<1$ and a compatible $\beta$ improve reward, convergence, or
stability over the base scheme, and whether aggressive selection must be
paired with a smaller $\beta$ to avoid excessively large per-sample shifts.

\paragraph{Scheme C: asymmetric hard selection.}
The matched $3\times3$ sweep tests whether positive samples require stricter
confidence filtering while negative samples benefit from full coverage.
Comparing B and C at the same $(\rho,\beta)$ isolates this asymmetry.  Better
and equally stable results from C would support keeping all negative samples
and filtering only borderline positive samples; better results from B would
favor symmetric filtering, while a best result at $\rho=1$ would indicate no
benefit from hard selection.  Sweeping $\beta$ also determines whether any
such conclusion is robust to the overall shift strength rather than an
artifact of the original $\beta=1$ setting.
